\begin{chapintro}
\kkref{sec:13-5}の主解析、カットオフ 0.5 ng/mL の 6 研究の要約点（感度 0.746、特異度 0.757）を、1000 人に PCT を測ったときの人数に直します。そのうえで、GRADE で結果の確実性を判定し、PRISMA-DTA に沿って報告を書きます。最後に、結果の段落の例文を示します。図 15.1 を描くコードと、検査前確率と検査後確率の図は、発展の印をつけて載せました。
\end{chapintro}

\section{1000 人あたりの数にする}
\label{sec:15-1}

感度 0.746 と特異度 0.757 という数字を聞いても、検査を使う場面で何が起きるかはすぐには浮かびません。同じ要約点を「1000 人の発熱患者に PCT を測ったら、何人が正しく判定され、何人が誤って判定されるか」に直すと、検査の役割に照らして結果を判断できるようになります。

人数に直すには有病割合を決める必要があります。組み入れた 10 研究の有病割合は、救急外来の研究で 11〜14\%、ICU の研究で 21〜23\% でした。そこで、研究の有病割合を見やすい数に丸めて、救急外来を想定した 10\%、ICU を想定した 25\% の 2 つの場面を置きます。計算は単純で、1000 人のうち菌血症の人数に感度を掛ければ真陽性、菌血症でない人数に特異度を掛ければ真陰性になり、残りがそれぞれ偽陰性と偽陽性です。括弧内は、感度と特異度の 95\% 信頼区間の端から計算した範囲です。1000 人あたりで示す考え方は、介入のメタ解析の結果の要約表（\kkref{sec:07-9}）と同じです。

\begin{kkcenter}
\kkfit{\begin{tabular}{lcc}
\toprule
1000 人あたり & 有病割合 10\%（救急外来） & 有病割合 25\%（ICU） \\
\midrule
菌血症の人 & 100 & 250 \\
\quad 真陽性 & 75（64〜83） & 187（159〜208） \\
\quad 偽陰性 & 25（17〜36） & 63（42〜91） \\
菌血症でない人 & 900 & 750 \\
\quad 真陰性 & 681（653〜707） & 568（544〜589） \\
\quad 偽陽性 & 219（193〜247） & 182（161〜206） \\
\midrule
陽性的中割合 & 0.25 & 0.51 \\
陰性的中割合 & 0.964 & 0.899 \\
\bottomrule
\end{tabular}}
\kkcap{表 15.1　カットオフ 0.5 ng/mL の 6 研究の要約点（感度 0.746、特異度 0.757）を 1000 人に用いたときの人数}
\end{kkcenter}

同じ人数を、1 人を 1 つの点にして並べると、誤りの多さを目で見られます。有病割合 10\% の 1000 人の場合を図 15.1 に示します。

\kkfig[.9\linewidth]{x2_icon1000.pdf}{図 15.1　有病割合 10\% の 1000 人を 1 人 1 点で並べた図（カットオフ 0.5 ng/mL の 6 研究の要約点、感度 0.746、特異度 0.757）}

点の数は、表 15.1 と同じ計算で求めます。\kkref{sec:13-4}の主解析の glmer() の結果 g（相関を 0 にしたモデル）から、感度と特異度を取り出します。

\begin{rcode}
se <- plogis(fixef(g)[["sens"]])         # 感度 0.746
sp <- plogis(fixef(g)[["spec"]])         # 特異度 0.757
tp <- round(100 * se); fn <- 100 - tp    # 菌血症の 100 人
tn <- round(900 * sp); fp <- 900 - tn    # 菌血症でない 900 人
c(TP = tp, FN = fn, FP = fp, TN = tn)
\end{rcode}

\begin{rout}
 TP  FN  FP  TN
 75  25 219 681
\end{rout}

点 1 つが 1 人で、濃い赤が真陽性（75 人）、薄い赤が偽陰性（25 人）、濃い青が偽陽性（219 人）、薄い青が真陰性（681 人）です。赤の系統の 100 人が菌血症の人、青の系統の 900 人が菌血症でない人にあたります。PCT が陽性になるのは濃い赤と濃い青を合わせた 294 人で、そのうち菌血症は 75 人です。濃い青の点の多さが、陽性的中割合が 0.25 にとどまる理由を示しています。見落とされる薄い赤の点は 25 個で、数は少ないものの、除外に使う場合はこの 25 人が問題になります。有病割合を 25\% にすると、赤の系統が 250 点に増え、青の系統が 750 点に減るので、同じ図でも濃い赤と濃い青の比が変わります。

陽性的中割合は、検査が陽性になった人のうち本当に菌血症である人の割合で、一般には陽性的中率とも呼ばれます。救急外来の場面では 75 ÷ 294 で 0.25 です。ICU の場面では有病割合が高いので、同じ感度と特異度でも陽性的中割合は 0.51 に上がります。\kkref{sec:09-3}で述べたとおり、感度と特異度が同じでも、的中割合は有病割合によって変わります。

\subsection{\hatten 図 15.1 を描くコード}
\label{sec:15-1-1}

rep() に人数を渡して、真陽性、偽陰性、偽陽性、真陰性の順に 1000 個のラベルを作り、50 列 20 行の格子に左上から順に置いています。

\begin{rcode}
lab <- rep(c("TP", "FN", "FP", "TN"), c(tp, fn, fp, tn))
col <- c(TP = "red", FN = "pink", FP = "blue", TN = "lightblue")
x <- rep(1:50, times = 20)     # 50 列
y <- rep(20:1, each = 50)      # 20 行
plot(x, y, pch = 19, col = col[lab], asp = 1,
     axes = FALSE, xlab = "", ylab = "")
\end{rcode}

\subsection{有病割合と的中割合}
\label{sec:15-1-2}

表 15.1 は有病割合 10\% と 25\% の 2 点だけを示しました。有病割合を連続的に動かしたときの陽性的中割合と陰性的中割合を、図 15.2 に示します。感度と特異度は主解析の要約点の値に固定しています。

\kkfig[.8\linewidth]{x2_ppv_prev.pdf}{図 15.2　有病割合と的中割合の関係（カットオフ 0.5 ng/mL の 6 研究の要約点、感度 0.746、特異度 0.757 に固定）。縦の点線は有病割合 10\% と 25\%}

的中割合は、1000 人の表の計算を式にしたものです。真陽性は感度 × 有病割合、偽陽性は（1 − 特異度）×（1 − 有病割合）に比例するので、陽性的中割合はその 2 つの比で決まります。

\begin{rcode}
prev <- seq(0.01, 0.6, by = 0.005)     # 有病割合
ppv <- se * prev /
  (se * prev + (1 - sp) * (1 - prev))  # 陽性的中割合
npv <- sp * (1 - prev) /
  (sp * (1 - prev) + (1 - se) * prev)  # 陰性的中割合
plot(100 * prev, ppv, type = "l", col = "red", ylim = c(0, 1),
     xlab = "有病割合（%）", ylab = "的中割合")
lines(100 * prev, npv, col = "blue")
abline(v = c(10, 25), lty = 3)
\end{rcode}

横軸は検査を受ける集団の有病割合、縦軸は的中割合で、赤が陽性的中割合、青が陰性的中割合です。有病割合が上がると、赤の曲線は右上がりに上昇し、青の曲線は下がります。有病割合 10\% では陽性的中割合が 0.25、陰性的中割合が 0.964、25\% では陽性的中割合が 0.51、陰性的中割合が 0.899 で、表 15.1 と一致します。同じ PCT でも、救急外来と ICU のように有病割合が違う場面では、陽性の意味が違います。

曲線は感度と特異度が有病割合によらず一定という前提で引いています。場面によって感度が変わる可能性もあるので、曲線は、検査を使う場面の有病割合を決めたうえで、その場面で推定した感度と特異度を使って描き直すのが原則です。

\subsection{\hatten 検査前確率と検査後確率}
\label{sec:15-1-3}

\kkref{sec:09-4}で、検査前の確率に尤度比を掛けると検査後の確率になることを、Kato 2010 の値で計算しました。要約点から求めた尤度比（陽性 3.07、陰性 0.34）を使い、検査前の確率ごとに検査後の確率を描いたものが図 15.3 です。

\kkfig[.65\linewidth]{x2_pretest.pdf}{図 15.3　検査前確率と検査後確率（カットオフ 0.5 ng/mL の 6 研究の要約点から求めた陽性尤度比 3.07、陰性尤度比 0.34）。点は検査前確率 10\% の場合}

\begin{rcode}
lrp <- se / (1 - sp)         # 陽性尤度比
lrn <- (1 - se) / sp         # 陰性尤度比
post <- function(p, lr) {    # 検査前確率 p から検査後確率を求める
  odds <- p / (1 - p) * lr
  odds / (1 + odds)
}
p0 <- seq(0.001, 0.999, length.out = 300)
plot(p0, post(p0, lrp), type = "l", col = "red",
     xlab = "検査前確率", ylab = "検査後確率")
lines(p0, post(p0, lrn), col = "blue")
abline(0, 1, lty = 3)        # 検査をしない場合
round(post(0.10, c(lrp, lrn)), 3)
\end{rcode}

\begin{rout}
[1] 0.254 0.036
\end{rout}

post() は、確率をオッズに直して尤度比を掛け、確率に戻す関数です。

横軸は、検査を行う前に、その患者が菌血症である見込みです。縦軸は検査の結果が分かった後の見込みです。点線の対角線は検査をしない場合で、検査前と検査後が同じ値になります。赤の曲線は PCT が陽性のとき、青の曲線は陰性のときの検査後確率です。検査前確率が 10\% のとき、陽性なら 0.254、陰性なら 0.036 になります。10\% は表 15.1 の救急外来の有病割合と同じなので、0.254 は陽性的中割合 0.25 に、0.036 は 1 − 陰性的中割合 0.964 に一致します。

赤と青の曲線が対角線から離れている幅が、この検査で見込みが動く大きさを表します。この検査の尤度比では、陽性でも検査前確率 10\% の人の見込みは 3 割に届かず、陰性でも 0.036 が残ります。検査前確率は患者ごとに違い、診察の所見や経過で変わるので、図 15.3 の曲線は、個々の患者の検査前確率を推定したうえで読みます。

\subsection{検査の役割から誤りの重さを決める}
\label{sec:15-1-4}

表 15.1 の誤りのどちらを重く見るかは、PCT をどの役割で使うかで決まります。

\begin{description}
\item[除外に使う場合] PCT が陰性なら血液培養や抗菌薬を見送る、という使い方です。問題になるのは偽陰性で、救急外来の場面では 1000 人あたり 25 人（17〜36 人）の菌血症を見落とします。陰性的中割合 0.964 は高く見えますが、見落とされた 25 人は治療が遅れる人たちです。この数を許容できるかは、臨床の判断です。
\item[確定に使う場合] PCT が陽性なら菌血症として扱う、という使い方です。問題になるのは偽陽性で、救急外来の場面では 1000 人あたり 219 人（193〜247 人）が不要な精査や抗菌薬の対象になります。
\item[追加の検査として使う場合] 診察とほかの検査で判断がつかない人にだけ PCT を加える、という使い方です。この場合、関わる患者の有病割合は表 15.1 のどちらとも違う可能性があり、その集団の精度データが必要になります。
\end{description}

表 15.1 の人数から言えるのは、PCT 単独では菌血症を除外することも確定することもできない、ということです。要約点が示す精度の水準と、\kkref{sec:13-6}で見た予測領域の広さの両方が、その判断の根拠になります。

\section{GRADE で確実性を判定する}
\label{sec:15-2}

GRADE は\kkref{sec:07-8}で介入のメタ解析に当てはめました。検査の精度にも同じ枠組みを使いますが、いくつか違いがあります。研究デザインとバイアスのリスク、非直接性の扱いは Schünemann らの GRADE guidelines 21 part 1、非一貫性、不精確さ、出版バイアスの扱いと結果の要約表の作り方は part 2 にまとまっています（Schünemann 2020a、2020b）。

\subsection{検査の GRADE で違うところ}
\label{sec:15-2-1}

1 つめは出発点です。検査の精度の研究は、対象を連続して登録し、全員に対象検査と参照基準を行う横断研究であれば、確実性「高」から始めます。介入で RCT から始めるのと同じ位置づけです。症例対照デザインのような研究は、バイアスのリスクの領域で下げます。

2 つめは、判定の単位です。感度に関わる真陽性と偽陰性、特異度に関わる真陰性と偽陽性について、それぞれ確実性を判定します。感度の研究は菌血症の人だけ、特異度の研究は菌血症でない人だけから推定されるので、ばらつきや不精確さの程度が両者で違うからです。

3 つめは、非直接性の考え方です。検査の精度は、患者にとっての結果（死亡、入院期間、抗菌薬の使用量）を直接示すものではありません。検査の結果に基づいて治療が変わり、その治療が患者の結果を変える、という一連の流れを前提にしています。

\subsection{検査の精度と患者の転帰}
\label{sec:15-2-2}

検査の精度が高いことは、その検査を使うと患者の転帰が良くなることを意味しません。検査の結果が診療を変え、その変更が転帰を改善して初めて、患者に利益が生じます。精度の研究が示すのは、検査が病気のある人とない人を見分けられるかどうかまでです。検査の結果を使った診療と使わない診療を比べ、死亡や入院期間のような患者の転帰を調べるのは、検査と治療を組にした無作為化試験の役割です。そうした試験がないレビューでは、精度の結果を患者の利益として読まず、結果の要約表に注記を付けます。

\subsection{5 つの領域}
\label{sec:15-2-3}

確実性を下げる領域は、介入と同じく 5 つあります。

\begin{description}
\item[バイアスのリスク] QUADAS-2（または QUADAS-3）の判定から決めます（判定の使い方は\kkref{sec:12-8}）。高リスクの研究が要約点の多くを占めていれば下げます。
\item[非直接性] 研究の対象、対象検査、参照基準が、レビューの問いとずれていないかを見ます。測定試薬が国内で使われているものと違う、血液培養のセット数が少なく参照基準が不完全である、といったずれがここに入ります。
\item[非一貫性] 研究間のばらつきが大きく、その原因が説明できないときに下げます。I² は検査のメタ解析では使いにくいので（\kkref{sec:10-7}）、研究ごとの感度と特異度の散らばりと、予測領域の広さで判断します。
\item[不精確さ] 信頼区間の幅を、1000 人あたりの人数で評価します。偽陰性が 17 人の場合と 36 人の場合とで、臨床の判断が変わるかどうかを考えます。
\item[出版バイアス] 診断精度の研究では、試験登録が普及しておらず、小規模研究効果を調べる検定の検出力も低いので、評価は難しくなります。介入の Egger 検定は使わず、Deeks の方法を使います（\kkref{sec:10-8}、Deeks 2005）。
\end{description}

\subsection{PCT の例での判定}
\label{sec:15-2-4}

主解析のカットオフ 0.5 ng/mL の 6 研究について、架空の PCT のレビューで判定した例を表 15.2 に示します。判定の理由は、レビューの著者が 1 つずつ書き残します。ここでの判定は説明のための架空のものです。

\begin{kkcenter}
\kkfit{\begin{tabular}{lll}
\toprule
領域 & 感度（真陽性と偽陰性） & 特異度（真陰性と偽陽性） \\
\midrule
バイアスのリスク & 1 段階下げる & 1 段階下げる \\
非直接性 & 下げない & 下げない \\
非一貫性 & 1 段階下げる & 下げない \\
不精確さ & 下げない & 下げない \\
出版バイアス & 下げない & 下げない \\
\midrule
確実性 & 低 & 中 \\
\bottomrule
\end{tabular}}
\kkcap{表 15.2　カットオフ 0.5 ng/mL の 6 研究の要約点についての GRADE の判定（架空の例）}
\end{kkcenter}

判定の根拠を領域ごとに書きます。

\begin{description}
\item[バイアスのリスク] 6 研究のうち 2 研究（Rossi 2019 と Qureshi 2018）が QUADAS-2 で高リスクでした。Rossi 2019 はデータを見てカットオフを決めており、Qureshi 2018 は一部の患者にしか参照基準を行っていません。この 2 研究を除いた 4 研究の要約点は感度 0.709（0.573〜0.815）、特異度 0.781（0.755〜0.806）で、感度は下がり、特異度は上がりました（\kkref{sec:14-5}）。感度、特異度とも 1 段階下げました。
\item[非直接性] 対象の場面（救急外来と ICU）、対象検査（PCT）、参照基準（血液培養）は、レビューの問いと合っているので、下げませんでした。患者の転帰への影響は精度の研究からは分からないので、結果の要約表に注記します。
\item[非一貫性] 6 研究の感度は 0.56 から 0.88 まで散らばり、研究間の標準偏差は logit 感度で 0.518 でした。カットオフはそろえてあり、場面とカットオフを調べても原因を一つに決められなかったので（\kkref{sec:14-2}）、感度は 1 段階下げました。特異度は 0.71 から 0.80 の範囲に収まり、logit 特異度の研究間の標準偏差は 0.136 と小さいので、下げませんでした。
\item[不精確さ] 6 研究で、感度の信頼区間は 0.638 から 0.831 と、特異度の 0.726 から 0.786 より広く開いています。有病割合 10\% の救急外来の場面で、偽陰性は 17〜36 人、偽陽性は 193〜247 人の範囲にあります。どちらの端でも「PCT 単独では除外も確定もできない」という判断は変わらないので、下げませんでした。
\item[出版バイアス] 6 研究では、小規模研究効果を調べる Deeks の検定で出版バイアスを確かめることが難しいと記載したうえで、未発表の研究がある証拠もないので、下げませんでした。
\end{description}

\begin{hattenbox}
Deeks の検定は、診断オッズ比に研究間の異質性があるときに検出力がとくに低く、有意でなくても出版バイアスがないとは言えません（Cochrane の DTA ハンドブック 9.5.3）。6 研究の診断オッズ比は 4.56 から 23.88 まで開いています。
\end{hattenbox}

この判定では、感度の確実性は「低」、特異度は「中」です。感度と特異度を別々に判定したために、2 つの確実性が分かれました。

表 15.1 と表 15.2 を合わせたものが、結果の要約表（Summary of findings 表）の中身になります。PCT の例なら、真陽性、偽陰性、真陰性、偽陽性の 4 行それぞれに、1000 人あたりの人数（有病割合 10\% と 25\%）、研究数と人数（6 研究、1,782 人で、菌血症 272 人と菌血症でない 1,510 人）、確実性を並べます。

\section{PRISMA-DTA に沿って報告する}
\label{sec:15-3}

診断精度のメタ解析の報告指針は PRISMA-DTA です（McInnes 2018）。介入のレビューの PRISMA を、診断精度のレビューに合わせて項目を変えたり加えたりしたもので、多くの雑誌が投稿時にチェックリストの提出を求めます。学会発表や論文の抄録については、抄録版の PRISMA-DTA for Abstracts があります（Cohen 2021）。\kkref{sec:08-1}で見た PRISMA 2020 と共通する項目が多く、検査のレビューで書き落としやすいのは、論文の部分ごとに次の点です。

\begin{description}
\item[題と抄録] 題で診断精度のレビューであることが分かるようにします。抄録には、対象検査、参照基準、標的疾患、要約の感度と特異度を書きます。
\item[序論] 対象検査を臨床のどの役割（置き換え、トリアージ、追加）で使うことを想定しているかを書きます。
\item[方法] 対象、対象検査、参照基準、標的疾患を定義します。カットオフの扱い（どのカットオフを主解析にするか、複数あるときにどれを取ったか）、バイアスのリスクの評価方法（QUADAS-2 など）、統計手法（二変量モデル、HSROC モデル、使ったソフトウェアとパッケージ）、ばらつきを調べる共変量を、事前に決めたものとして書きます。
\item[結果] 研究ごとの 2×2 表、感度と特異度と信頼区間を表かフォレストプロットで示します。SROC 曲線、要約点、信頼領域、予測領域を図で示します。QUADAS-2 の判定を研究ごと、領域ごとに示します。
\item[考察] 結果を臨床の言葉（1000 人あたりの人数）で述べ、ばらつきとバイアスのリスクがどの程度結論を弱めるかを書きます。
\end{description}

研究ごとの 2×2 表は、読者がメタ解析をやり直すための最低限のデータです。本文の表に収まらなければ、補足資料に載せます。解析に使った R のコードも、補足資料かリポジトリで公開しておくと、読者が結果を確かめられます。

\section{結果の段落の例文}
\label{sec:15-4}

架空の PCT のレビューについて、結果の段落を書くと次のようになります。数値は本書の例題のものです。

\begin{KKfbox9ao}[例文]
検索で得た 2,315 件から、10 研究（3,626 人、うち菌血症 556 人）を組み入れた。6 研究が救急外来、4 研究が ICU の患者を対象とし、PCT のカットオフは 0.25、0.5、1 ng/mL であった。カットオフが混ざるため、主解析は研究数の最も多いカットオフ 0.5 ng/mL の 6 研究（1,782 人、うち菌血症 272 人）とし、二項分布の二変量モデルで要約した。研究間の相関が −1 の境界に張り付いたため、相関を 0 にしたモデルに単純化した。要約の感度は 0.746（95\% 信頼区間 0.638〜0.831）、特異度は 0.757（0.726〜0.786）であった。6 研究の感度は 0.56 から 0.88、特異度は 0.71 から 0.80 の範囲にあった。全 10 研究は SROC 曲線とカットオフの共変量の解析に使い、カットオフが高いほど偽陽性割合が低かった（logit 偽陽性割合の係数 −0.77、p = 0.006）。バイアスのリスクが高い 2 研究（Rossi 2019、Qureshi 2018）を除いた 4 研究では、感度 0.709（0.573〜0.815）、特異度 0.781（0.755〜0.806）で、特異度の研究間の標準偏差は 0 と推定された。有病割合 10\% の 1000 人に用いると、菌血症 100 人のうち 25 人（17〜36 人）が偽陰性、菌血症でない 900 人のうち 219 人（193〜247 人）が偽陽性になると推定された。確実性は、感度を「低」、特異度を「中」と判定した。
\end{KKfbox9ao}

例文はである調で書きました。論文の本文は、雑誌の慣習に合わせて文体を選びます。数値の示し方は、主解析の対象と要約点、ばらつきの範囲、共変量の結果、感度分析、1000 人あたりの人数、確実性の順にすると、読者が結論の強さを一度に判断できます。

\section{この章のまとめ}
\label{sec:15-5}

要約点は、検査を使う場面の有病割合を置いて 1000 人あたりの人数に直し、検査の役割から見て重い誤り（除外なら偽陰性、確定なら偽陽性）の数で判断します。精度が高いことと患者の転帰が良くなることは別の問いで、後者は検査と治療を組にした試験で調べます。確実性は GRADE で、感度と特異度を別々に判定し、その理由を書き残します。報告は PRISMA-DTA に沿い、研究ごとの 2×2 表と解析コードを公開します。第 2 部はここまでで、問いを立ててから報告するまでの工程を、検査のメタ解析で一通りたどりました。この先を学ぶには、Cochrane の DTA ハンドブックの第 9 章（考え方）と第 10 章（実行）が最も詳しい資料です（付録 C）。
